
\section{هفته دوم}
\subsection{دوشنبه ۲۴ دی}
بله امروز دیگه قرار نیست از من غر بشنوید. میریم سراغ بررسی \cl ها.
اول بریم سراغ مرج شده ها.

\subsubsection{مرج شده ها}
{\LARGE go/FIPS}

اولین کامیت مرج شده ای که به چشمم خورد
\LTRfootnote{\url{https://go-review.googlesource.com/c/go/+/639775}}
راجع به اضافه کردن تست داخل زیرپوشه \lr{src/crypto/internal/} هست. اینکه چی هستن و چیکار میکنن مربوط FIPS
\LTRfootnote{\lr{Federal Information Processing Standards}}
هستند و یک سری استاندارد امنیتی هست که یکی از استاندارد های شناخته شدش \lr{FIPS 140} هست
\LTRfootnote{\url{https://en.wikipedia.org/wiki/FIPS_140}}
که مقوله جالبیه ولی الان فرصت ندارم زیاد راجع بهش بخونم. این داستان و \cl از یک ایشو که شبیه به یک پروپوزال هم هست شروع شده
\LTRfootnote{\url{https://github.com/golang/go/issues/69642}}
و همون رو آپدیت کرده. وقتی ایشو رو بررسی کنید اطلاعات جالبی پیدا میکنید. اول که مشکلاتی رو توضیح داده. بعد از اون میبینید که یه عالمه
کامیت داشته که همین ایشو رو آپدیت میکنن. همین نشون میده ما میتونیم یه سری \cl ارسال کنیم که بخشی از نیاز یک ایشو رو صرفا انجام بده.

از اونجایی که متن کامیت مهم هست میذارمش:

\begin{latin}
\begin{minted}[fontsize=\footnotesize,breaklines]{text}
crypto/internal/fips140test: add hmac DRBG ACVP tests

Adds ACVP test coverage for the hmacDRBG algorithm based on the NIST
spec:
  https://pages.nist.gov/ACVP/draft-vassilev-acvp-drbg.html#section-7.2

The HMAC DRBG algorithm in our fips module is a minimal implementation
tailored for use for generating ECDSA nonces and so lives in
crypto/internal/fips140/ecdsa.

In order to be testable by crypto/internal/fips140test this changeset
exports a ecdsa.TestingOnlyNewDrbg() constructor to support the ACVP use-case.
All FIPS-compatible SHA2 and SHA3 digests are tested.

The ACVP capability registration is customized to match the limited
capabilities of our ecdsa-focused impl. Most notably:

  * reseedImplemented is false - we expect this impl to be invoked
    only once or twice per instantiation and do not support explicit
    reseeding.
  * predResistanceEnabled is false - this requires reseeding.
  * Per mode:
    * derFuncEnabled is always false - this is only used by ctrDRBG.
    * additionalInputLen is 0 for all modes - this is only used with
      preResistanceEnabled.

The other capability values are chosen based on Table 4:
  https://pages.nist.gov/ACVP/draft-vassilev-acvp-drbg.html#section-7.4

Updates #69642

Change-Id: Ia58979d691f912e2ed739a05efb719f580fbbf89
Reviewed-on: https://go-review.googlesource.com/c/go/+/639775
Reviewed-by: Michael Pratt <mpratt@google.com>
LUCI-TryBot-Result: Go LUCI <golang-scoped@luci-project-accounts.iam.gserviceaccount.com>
Auto-Submit: Filippo Valsorda <filippo@golang.org>
Reviewed-by: Dmitri Shuralyov <dmitshur@google.com>
Reviewed-by: Filippo Valsorda <filippo@golang.org>
\end{minted}
\end{latin}

نکات خوبی از همین کامیت مسیج میشه یاد گرفت. خط اول رو ببینید و البته لینک هایی که در ادامه گذاشته شدن. پایین تر هم مهم ترین نکته از نظر من
همون \lr{Updates \#issue} هست. به‌نظرم خیلی مسئله مهمی هست. ریپو این پروژه هم که مشخصا \lr{go} بود.

یک نکته دیگه ای که از این میشه متوجه شد، پروژه های اوپن سورس بخش های خیلی متنوعی دارن. برای کمک توشون باید یه تیکه خاص رو خیلی خوب
عمیق بشیم و بعد بریم یاری برسونیم. مثلا اگه بخوام مشابه کار همین اقا رو انجام بدم باید بشینم استاندارد های مربوطه رو بخونم و کد بزنم.

موضوع مهم دیگه ای که میشه بررسی کرد کامنت های \lr{Gerrit} هستن. از پایین کامنت ها رو دیدم واسه همین \cl. اول خب دوتا کامنت بود
که یک سری مشکلات جالب و ساده رو ریپورت کرده بودن و فیکس شد. لحن جفتشون هم خیلی دوستانه و خوب بود.

\begin{latin}
\textbf{Filippo Valsorda(Maintainer):}

TestingOnlyNewDRBG

Acronyms go all in the same case as the first letter.

\textbf{Daniel McCarney(Developer):}

I swear I will remember this eventually :)

Fixed in patchset 5. Thanks
\end{latin}

اون آقای \lr{Filippo} رو دیدم خیلی آدم خفنی بود. ۱۰ سال هست که طرف داخل تیم گو هست و کلی کار جالب درمورد رمزنگاری کرده.
این مثلا منتور خوبی هست. جای نگرانی نبود تا رسیدم به کامنت بعدی. یک چیزایی راجع به چرخه رلیز گو
\LTRfootnote{\url{https://go.dev/wiki/Go-Release-Cycle}}
گفته بود که نخوندمش و الان میرم مطالعش می‌کنم. درواقع چرخه انتشار خود مخزن گو خیلی متفاوت هست با پکیج های دیگه اون.
اینطور نیست که هر زمان سال که دوست داشتید تست ارسال کنید یا فیچر پیاده کنید. فعلا هم داخل دوره ای هستیم که فقط مستندسازی و فیکس کردن باگ
های حساس پذیرفته میشه. حالا باز جلوتر راجع به این \lr{Go Release Cycle} صحبت می‌کنیم. البته که تو همون کامنت \lr{Filippo} توجیه کرد که نخیر این کامیت
لازم هست و بدردمون میخوره. زود باش مرجش کن. و خب این شد که با پارتی بازی این کامیت مرج شد. البته ممکنه اشتباه متوجه شده باشم داستان رو.

کاربری که این کامیت رو ارسال کرده هم توسعه دهنده خیلی باتجربه ای هست. از ۲۰۱۰ اکانت گیتهاب داره و البته ۱۳ تا کامیت مرج شده هم داره.
به کارنامش که نگاه کنید میبینید مطالعه روی \lr{FIPS} و ساپورتش برای گو تو برنامه هاش هست. رسما واسه مشارکت موفق باید یک موضوع رو عمیق
شد و مدت ها روش کار کرد. نمیشه همینطور هر ماه رفت موضوع جدید. پس همین بهم میگه قبل از شروع جدی باید یک چیزی رو پیدا کنم تا توش
مدت ها مشارکت کنم. این مشارکت کننده اونقدر خفن نبود ولی میشه حتی بهش ایمیل زد و برای منتور شدن ازش سوال پرسید. این هم پلن جوابیه.
با خفنا ایمیل بازی کنی و ازشون بخوای راهنماییت کنن.

{\LARGE چرخه انتشار مخزن گو}

مستند چرخه انتشار مخزن گو یا \lr{Go Release Cycle}
\LTRfootnote{\url{https://go.dev/wiki/Go-Release-Cycle}}
هست. من نمیخوام الان این مستند رو براتون ترجمه کنم. برید مطالعه کنید ولی چند نکته مهم رو میگم. اول اینکه گو هر شیش ماه رلیز داره
و برای هر کدوم از این رلیز هاش حدودا ۵ ماه مخزن باز هست و اجازه هرطور مشارکتی رو میده و بعد ۲/۳ ماه فریز میشه. برای مثال همین الان
برای رلیز زمستونی توی فریز هستیم و دست بنده برای مشارکت داخل مخزن گو و زیرمخزن هایی که وابستگی زیاد دارن یکم بسته میشه.
زمانی که داخل فریز هستیم، تنها کامیت‌هایی که بهتره بگم \cl بهشون پذیرفته می‌شه باگ فیکس، تست و تغییر و اضافه کردن مستندات هستن.
درواقع بعد از پروسه توسعه یک رلیز، یک مدت تیم گو میخوان از کیفیت رلیز مطمعن بشن و تو اون مدت فقط پروژه رو تست میکنن و باگ های مهم رو
بررسی میکنن.

حالا یک نکته کلیدی که اول مستند نوشته. آقا غیر از ایشو ها که میشه داخلشون اول مشکل رو شرح بدیم و بعد \cl  بفرستیم یه جای دیگه هم هست
داخل \lr{google group} ها. داخل همین صفحه چرخه انتشار برید آدرس گوگل گروپ هست. پیام و سوال یا هرچیزی رو هم میتونید به ایمیل
\href{conduct@golang.org}{conduct@golang.org}
ارسال کنید و داخل همون گوگل گروپ لیست میشه. این لیست یکی از فایده هاش این هست که نمیخواد خودتون حساب کنید الان کجای \lr{Release Cycle}
هستیم. البته که اصلا نمیتونید... به‌خاطر اینکه الزاما تاریخ ها مشخص نیستن و خیلی وقتا تقریب زده میشن. اینکه کی فریز داریم و کی رلیز داریم و این حرفا
رو از همین گوگل گروپ میتونید دنبال کنید.
البته از حق نگذریم گوگل گروپ خیلی خلوت هست اما حدس میزنم راحت جواب بگیره آدم ازش. بعد هم که ایشو میشه زد راحت. کلا مستقیم کد فرستادن
زیاد باب میل نیست. مگر اینکه روی ایشو خاصی که قبلا کسی کار میکرده آپدیت یا فیکس انجام بدیم.

نکته نهایی اینکه برای ابتدای کار ممکنه یکم مخزن گو از این بابت که الان تایم فریز هست سخت تر هم باشه. شمایی که قصد مشارکت دارید هم ممکنه
وضعیت من رو داشته باشید. شروع از مخازن فرعی شاید راحت تر باشه. هرچند زیاد رضایت بخش نیست از نظر من.


{\LARGE cmd/go}

درفاصله نوشتن همین گزارش ها 18 کامیت دیگه مرج شدن. این نشون میده شانس مرج شدن خیلیه چون خیلی پیگیر هستن تیم توسعه گو.
از این بابت خوشحالم. یعنی هر روز جواب میگیرم. هر روز میتونم تلاش کنم تا مشارکتم ثبت بشه.

این یکی \cl خیلی نکته ای نداشت. ;کامیت و کامنت ها رو میذارم براتون.
\begin{latin}
\begin{minted}[fontsize=\footnotesize,breaklines]{text}
cmd/go: check go version when parsing go.mod fails

Fixes #70979

Change-Id: I6597fe178eed34702eea6cba4eec5174c9203458
Reviewed-on: https://go-review.googlesource.com/c/go/+/639115
LUCI-TryBot-Result: Go LUCI <golang-scoped@luci-project-accounts.iam.gserviceaccount.com>
Reviewed-by: Michael Matloob <matloob@golang.org>
Reviewed-by: Michael Knyszek <mknyszek@google.com>

--------------
comments
{\textbf Michael Knyszek}

conservatively added the wait-release hashtag, but if this is safe to land for Go 1.24, feel free to remove it and submit. thanks.

{\textbf Michael Matloob}

I think this is safe. I'll remove the hashtag and submit. Thanks!

\end{minted}
\end{latin}

یک نکته جالب بگم که \cl ها بر اساس حجمشون و تعداد خط کدی که کم و زیاد شده برچسب سایز دارن. این سایز ها هم مثل لباس هست.
\lr{S, M, L, XS, XL}


\subsection{سه‌شنبه ۲۵ دی}
امیدوارم همتون همیشه سلامت باشید. بیماری کم کارم کرده و کم کاری اذیتم می‌کنه. بریم سراغ کار مون.
\subsubsection{اوپن ها}
خب تعداد اوپن ها خیلی زیاد هست. مثلا تو پروژه گو برای امروز یک مرج داشتیم و الباقی همشون اوپن بودن و رد شده که طبیعی هم هست.
هرچقدر به قبل تر بریم تعداد مرج ها بیشتر میشن. به تاریخ های خیلی قبل تر که برید و مثلا ۲۰۰۰ \cl به قبل
\LTRfootnote{\url{https://go-review.googlesource.com/q/project:go,2000}}
رو بررسی کنید میبینید بهرحال یا رد شده یا مرج شده. پس جای خوشحالی هست که یک روزی بالاخره هر \cl ای پاسخ داده میشه
و به مسیر درست هدایت میشه.
البته اینکه یک \cl برای مثال یکی دو ماه باز بمونه و بلاتکلیف باشه هم غیر ممکن نیست.

بیایید در ابتدا بررسی کنیم کدوم \cl ها بعد مدت زیادی بررسی میشن و درواقع رها میشن.

اولین چیزی که به چشمم خورد این کامیت بود:
\url{https://go-review.googlesource.com/c/go/+/602015}
خب سایزش \lr{XS} هست و ۱۵۰۰ خط کد به کد گو اضافه کرده. به ایشو خاصی هم اشاره میکنه و متن کامیتش هم اوکی هست.
کلا این \cl ماه ها معطل شده و خودش هم زیاد براش مهم نبوده گویا. ولی کلا دیدم ۲ تا مرج شده داره ۱۰ تای دیگه رو هوا هستن.
دوتا مرج شدش هم مال مخزن \lr{benchmark} هست. من علتی که به ذهنم رسید انقدر دیر و بد کامیت هاش مرج شدن فقط بزرگ بودن کامیت بود.
ایده دیگه ای ندارم. باید بیشتر کامیت بررسی کنم ببینم بقیه چیه قضیه شون.

\subsection{چهارشنبه ۲۶ دی}
بری سر کامیت های اوپن از قدیم. این کامیت که واقعا باید باز باشه. دوتا کامنت داره که جواب نداده و \lr{resolve} نکرده.
و متن کامیت رو ببینید:

\begin{latin}
\url{https://go-review.googlesource.com/c/go/+/425875}

\begin{minted}{text}
Commit message:
---------
text/template: add "return" action

For #53015

Change-Id: I147cf4abec84d72a57436678b08c63c1de330e26
\end{minted}
\end{latin}

این دوستمون که واقعا حقش بوده باز باشه \cl اش. تازه الان کانفلیکت هم داره.

همین الان یک نکته. یک سری از \cl ها وقتی قدیمی می‌شن کلا مرج کانفلیکت پیدا میکنن تو فیلد های تغییرات \lr{LTP} یا \lr{Legacy TryBot Pass}
که باید پاس بشه. اگر -۱ بگیره یعنی مرج کانفلیکت داریم و داستانه. پس خیلی از \cl های اوپن مشکل کانفلیکت دارن و صرفا کامنت هاشون رو جواب ندادن
و پیگیر نبودن.

مورد بعدی که دارم بررسی می‌کنم

\begin{latin}
\url{https://go-review.googlesource.com/c/go/+/600875}

\begin{minted}[fontsize=\footnotesize,breaklines]{text}
Commit message:
---------
runtime/cgo: Better Handle concurrency

This change seeks to dramatically improve the performance of Handles in concurrent and non-concurrent settings. It does this through the use of
atomics, slices, and locks.

On my machine, I see 3-400% improvement over master, and a +/-20% increase over my previous commit.

How it works:

1. When creating new handles, we check a sync.Pool for any recent released indices
2. If we have one, we use it, otherwise we start from the beginning of the slice -- to keep the slice compact and avoid growing
3. If we reach the end of the slice, grow the slice by one element
4. Otherwise, try and set the slice to the new value or increment our index by one
5. Getting the value is somewhat trivial
6. And finally, when we delete a handle, we release the index back into the sync.Pool to be reused

Note that the sync.Pool is rather controversial as this would possibly be a BC break (previously, handles were monotonic). In this case, the
memory efficiency may be worth it. However, if this isn't a good idea, we can sacrifice a bit of performance/complexity for monotonicity without
giving up memory efficiency.

Here are the benchmark results on my machine:
goos: linux
goarch: amd64
pkg: runtime/cgo
cpu: 11th Gen Intel(R) Core(TM) i7-11800H @ 2.30GHz
                         │  master.txt  │              worst.txt               │              best.txt               │
                         │    sec/op    │    sec/op     vs base                │   sec/op     vs base                │
Handle/non-concurrent-16   469.65n ± 3%   139.45n ± 6%  -70.31% (p=0.000 n=10)   98.77n ± 3%  -78.97% (p=0.000 n=10)
Handle/concurrent-16        581.2n ± 3%    248.5n ± 3%  -57.24% (p=0.000 n=10)   104.1n ± 6%  -82.08% (p=0.000 n=10)
geomean                     522.5n         186.2n       -64.37%                  101.4n       -80.59%

best.txt: patch as-is (~O(1) freelist lookup)
worst.txt: adding to sync.Pool is commented out during deletes, forcing an ~O(n) freelist lookup

Change-Id: I0267b96807bf859967f3fbeac8081ea991da0713
\end{minted}
\end{latin}

این بابا روی پرفرمنس یک تیکه کد کار کرده و خیلی اذیتش کردن. تهش هم مرج نشده هنوز. ولی ایشالله اقا خیلی که پیگیری کنه مرج میشه.
یک کامنت خیلی آموزنده پیدا کردم میذارم اینجا. 
\begin{latin}
\textbf{Ian Lance Taylor:}

Thanks. The preferred way to measure performance is run the benchmarks several times with, for example, -test.count=10. Run them with and without your change. Pass the old and new output to the golang.org/x/perf/cmd/benchstat program. Put the benchstat output in the change description. Thanks.
\end{latin}
راستی این مشارکت کننده اولین مشارکتش بوده. شاید طبیعی هست که انقدر اذیتش کردن.

البته کلا بعضی از \cl ها رو میبینم نمیدونم چرا مرج نشدن. شاید باید کامنت بذاریم زیرشون و جویا بشیم که مرج کنید لطفا این بندگان خدارا.
این هم ایده خوبی برای فعال بودن در کامیونیتی هست.حتما یک پست داخل گوگل گروپ می‌نویسم بپرسم اوکیه کامیت های قدیمی رو کامنت بذارم
و مثلا \lr{maintainer} ها رو \lr{CC} کنم.

\subsubsection{ریجکت شده ها}
بررسی \lr{cl} های رد شده نیازی به بررسی زمانشون نداره. بریم اینا رو ببینیم.

یک سری \lr{cl} های رد شده اصلا مال \lr{GerritBot} هست. 
اینجا میبینیم که اصلا \lr{GerritBot} کامیت اشتباه کرده و برای همین رد شده.
\begin{latin}
\url{https://go-review.googlesource.com/c/go/+/638935}

\begin{minted}{diff}
diff --git a/src/database/sql/convert_test.go b/src/database/sql/convert_test.go
index 1b2e61c..7cb8f317 100644
--- a/src/database/sql/convert_test.go
+++ b/src/database/sql/convert_test.go
@@ -32,7 +32,7 @@
    wantint    int64
    wantuint   uint64
    wantstr    string
-   wantbytes  []byte
+   wantbytes  []byte[just new white space]
    wantraw    RawBytes
    wantf32    float32
    wantf64    float64
@@ -67,7 +67,7 @@
    return []conversionTest{
        // Exact conversions (destination pointer type matches source type)
        {s: "foo", d: &scanstr, wantstr: "foo"},
-       {s: 123, d: &scanint, wantint: 123},
+       {s: 123, d: &scanint, wantint: 123},[just new white space]
        {s: someTime, d: &scantime, wanttime: someTime},
 
        // To strings
\end{minted}
\end{latin}
انتهای خط مشخص کردم که صرفا یک اسپیس غیرضروری که اتفاقا باید پاک بشه اضافه کرده بود و برای همین رد شده. البته اون متن نبود و مشخصا
یک اسپیس خالی بود.

البته که \lr{GerritBot} صرفا \lr{Pull Request} خاصی داخل گیتهاب رو ارسال کرده این طرف! و خب این هم مشکل طرف بوده که فرستاده گیتهاب.
پس جاهایی که میبینیم \lr{Owner} تغییرات \lr{GerritBot} هست همه از گیتهاب میان و داخل کامیت با تگ \lr{GitHub-Pull-Request} هم مشخص شده
کدوم \lr{PR} بوده.

یک آمار کلی گرفتم. ده درصد \cl های مرج شده از گیتهاب اومدن و تقریبا بیست درصد از \cl های رد شده هم مربوط به گیتهاب هستن.
از گیتهاب \lr{PR} نفرستید دوستان زیاد باحال نیست. \lr{Workflow} پروژه استفاده از \lr{Gerrit} هست و حالا گیتهاب رو هم ساپورت کردن خوشحال باشیم
ماهایی که تازه‌کار هستیم.

ببینیم الباقی چیکار میکنن. یک \cl که باز هم از گیتهاب میاد
\LTRfootnote{\url{https://go-review.googlesource.com/c/go/+/391835}}
مربوط به این هست که فیچر اضافه کرده داخل \lr{exec.Command} که بیاد یه سری تغییر ایجاد کنه که برید ببینید. نهایتا تو کامنت ها اینطور بودن که
کجا بدرد میخوره این؟ خب پس نمیتونیم بگیم آره من دوست دارم عملکرد فولان طور باشه. حتما باید یک جایی کاربرد معرفی کنیم ازش.

این دوستمون
\LTRfootnote{\url{https://go-review.googlesource.com/c/go/+/639956}}
سعی کرده یک جا به جای استفاده از \lr{bool} از \lr{struct} استفاده کنه و مثلا یک مقدار در مموری صرف جویی کنه. و خب بدون هیچ کامنتی و توضیحی
تغییراتش رد شده. البته در نهایت داخل لاگ ها علت رد شدن نوشته میشه همیشه که این رو تازه دیدم. بذارید کامنت طرف رو بذارم براتون. بله قطعا
در مقابل سادگی اون یه ذره حافظه بی ارزش هست. این هم آموزنده بود. پس سادگی برای تیم گو میتونه خیلی مهم تر از این حرفا باشه.

\begin{latin}
\textbf{Jakub Ciolek:}

Abandoned

The original is simpler, not worth the churn.
\end{latin}

یک \cl دیگه دیدم که خود توسعه دهنده بسته بود.
\LTRfootnote{\url{https://go-review.googlesource.com/c/go/+/624775}}
خب جالبه توضیحی هم نداده. البته بعد از یک ماه بسته بود... خب شاید بهش بر خورده. کی میدونه؟ یک سری از \lr{cl} ها رو کلا خود توسعه دهنده ها
تشخیص میدن بدرد نمیخورن و خودشون می‌بندن. حتی دیدم بعضا داخل خود تیم گو بودن و خودشون هم این کار رو انجام داده بودن و رد کرده بودن.
البته نکته ای که واقعا هست اینه که لیست تغییرات پر از مرج شده ها و اوپن ها هستن. واقعا \cl هایی که رد میشن خیلی کم هستن و همین خوشحال کنندست.

با بررسی پروژه های دیگه متوجه شدم خب کم و بیش داخلشون اکتیو هست. اما احتمالا برای هرکدوم \cl  ارسال بشه بررسی میکنن. در ابتدا میخوام
همون مشارکت های ساده رو روی پروژه های ساده تر این تیم داشته باشم تا دستم بیاد درست اخلاقیاتشون رو.

\subsection{تسک های جدید برای اولین مشارکت!}

\begin{itemize}
\item
بررسی مخازن با دستور \lr{go vet} اون هم با آخرین ورژن گو
\item
فیکس کردن مشکل و پرکتیس ها و مطالعه راجع بهشون
\item
ایمیل زدن داخل گو گروپ و پرسجو راجع به اینکه این مشکل رو باید ایشو زد؟ بعد از هماهنگی اینجا \cl ارسال کرد یا چی؟
\item
ادامه گفت‌وگوها در \lr{Gerrit}
\item
مرج شدن و شادی
\end{itemize}

بدون صحبت بریم سراغ کار!

\subsubsection{بررسی مخازن با دستور \lr{go vet} اون هم با آخرین ورژن گو}
با اسکریپت زیر شروع کردم به ثبت تمام مشکلات جزئی. بعد چند مرج اولم این فایل رو ایشو میکنم و میگم سوتونا بریزید همه باهم فیکس کنیم.

\begin{latin}
\begin{minted}{bash}
GO=/src/go/bin/go
for d in $(ls -d */); do
    cd $d
    echo ============== $d ============== >> ../vets.txt
    $GO vet  -all ./... 2>> ../vets.txt
    cd -
done
\end{minted}
\end{latin}

از بروز ترین کامپایلر گو و درواقع آخرین نسخه دارم استفاده می‌کنم.

از این بین من مشکلات مخزن \lr{benchmark} به چشمم خورد. کم و جالب هست.

\begin{latin}
\begin{minted}[fontsize=\footnotesize,breaklines]{text}
# golang.org/x/benchmarks/driver
driver/driver.go:375:2: unreachable code
driver/driver.go:376:2: unreachable code
# golang.org/x/benchmarks/third_party/biogo-examples/igor/igor
third_party/biogo-examples/igor/igor/cluster.go:51:9: github.com/biogo/store/interval.IntRange struct literal uses unkeyed fields
# golang.org/x/benchmarks/sweet/benchmarks/biogo-igor
# [golang.org/x/benchmarks/sweet/benchmarks/biogo-igor]
sweet/benchmarks/biogo-igor/igor.go:68:30: golang.org/x/benchmarks/third_party/biogo-examples/igor/igor.GroupConfig struct literal uses unkeyed fields
\end{minted}
\end{latin}

یک یا دوتا از همین موارد حل بشن هم خوب هست برای شروع. رفع این خطا ها بسیار آسان است و اگر مشارکت های کوچک به همین اندازه ساده باشه تعداد زیادی انجام میدم.
اولین \cl من
\LTRfootnote{\url{https://go-review.googlesource.com/c/benchmarks/+/642895}}
ارسال شد و الان احساس هیجان زیادی دارم. تو بخش \lr{Change Log} اولین پیامم رو دریافت کردم که میذارم براتون.

\begin{latin}
\begin{minted}{text}
Congratulations on opening your first change. Thank you for your contribution!

Next steps:
A maintainer will review your change and provide feedback. See
https://go.dev/doc/contribute#review for more info and tips to get your
patch through code review.

Most changes in the Go project go through a few rounds of revision. This can be
surprising to people new to the project. The careful, iterative review process
is our way of helping mentor contributors and ensuring that their contributions
have a lasting impact.
\end{minted}
\end{latin}

چشم آقای گوفر مهربون. البته یکم من سرسری ارسال کردم. بهتر بود جزئیات بیشتری راجع به تغییراتی که دادم میدادم.
 اگر خواستن، عوض میکنم حالا و در آینده بیشتر این کار رو انجام میدم.
فعلا نباید انتظار داشته باشم در لحظه اتفاقی بیوفته و میتونم تا فردا صبر کنم. اینطور تغییرات حداقل چند ساعت طول میکشه تا بررسی و مرج بشن.

یک مورد خیلی باحال. وقتی مثل من یک فرق غریبه هستید برای توسعه دهنده های گو، طول میکشه تا بررسی کنن \cl ات رو.
اما کسایی داخل خود تیم گو هستن و ایمیلشون مال گوگل یا گولنگ هست خیلی زود بررسی میشن و حتی مرج میشن.
البته این فرایند بسیار منطقی و طبیعی است. میرم کتاب هایی که نخوندم رو شروع میکنم به مطالعه تا کدم بررسی بشه.

\subsection{جمعه ۲۸ دی}
بعد از ۲۷ ساعت \cl من مرج شد و الان داخل مخزن benchmarks مشارکت دارم. بتونم چند تا کامیت دیگه روی همین مخزن داشته باشم
بدک نیست. میام بالاتر. چون تعداد مشارکت کننده هاش کمتر هست راحت میشه رفت تو چارت بهترین توسعه دهنده هاش.
البته که میتونه پروژه بی مصرف و بدرد نخوری به حساب بیاد.

فعلا میخوام چند تا \cl دیگه از همین سبک ارسال کنم. چطوره هرچی \lr{go vet} ارور میده رو ساکت کنیم؟ لیست \cl هام رو میذارم پایین.
\begin{latin}
\url{https://go-review.googlesource.com/c/arch/+/642738}
\end{latin}
این \cl بعد از یکونیم ساعت به درستی مرج شد و خیلی سریع و عالی پیش رفت. مسیر فیکس کردن وارنینگ ها خیلی راحته کلا.
انگیزه و توانایی فیکس کردن باقیشون رو دارم. اما تنها امیدم به این هست که یکی این متن ها رو بخونه و بتونه یک مشارکت داشته باشه
و راهش رو شروع کنه. خب الگوریتم خیلی راحته. برای تمام پروژه‌ها \lr{go vet -all ./...} رو اجرا کن، برو داخل مخزنی که میخوای و
فیکس کن همه چیز رو بعد هم بفرست کد رو و قطعا مرج میشه.

یک ایشو
\LTRfootnote{\url{https://github.com/golang/go/issues/71306}}
درست کردم و داخلش ریپورت \lr{go vet} رو برای خود مخزن گو گذاشتم.


\begin{latin}
\begin{minted}{bash}
go vet -all ./...

# crypto/internal/fips140test
# [crypto/internal/fips140test]
crypto/internal/fips140test/check_test.go:89:18: possible misuse of unsafe.Pointer
crypto/internal/fips140test/check_test.go:111:23: possible misuse of unsafe.Pointer
crypto/internal/fips140test/check_test.go:117:19: possible misuse of unsafe.Pointer
# internal/abi
# [internal/abi]
internal/abi/escape.go:21:9: possible misuse of unsafe.Pointer
# internal/runtime/atomic_test
# [internal/runtime/atomic_test]
internal/runtime/atomic/atomic_test.go:101:7: possible misuse of unsafe.Pointer
# reflect
# [reflect]
reflect/type.go:2245:35: possible misuse of unsafe.Pointer
# reflect_test
# [reflect_test]
reflect/all_test.go:8632:18: possible misuse of unsafe.Pointer
# runtime
# [runtime]
runtime/alg.go:39:22: possible misuse of unsafe.Pointer
runtime/arena.go:487:10: possible misuse of unsafe.Pointer
runtime/arena.go:492:10: possible misuse of unsafe.Pointer
runtime/arena.go:768:7: possible misuse of unsafe.Pointer
runtime/arena.go:780:14: possible misuse of unsafe.Pointer
runtime/arena.go:785:14: possible misuse of unsafe.Pointer
runtime/arena.go:792:14: possible misuse of unsafe.Pointer
\end{minted}
\end{latin}

بخشیش رو اینجا گذاشتم. خب تنها اروری که وجود داره \lr{possible misuse of unsafe.Pointer} هست که گزارش دادم. در پاسخ هم ایشو من
به عنوان \lr{duplicated} بسته شد. اطلاعات رو از داخل خود ایشو بررسی کنید. ولی توضیحش رو میدم. درواقع این وارنینگ \lr{false positive} هست.
هیچ مشکلی با کد نیست ولی بیخودی ارور میده. بعضی اوقات \lr{vet} چرتوپرت میگه. و این مشکلیه که داخلش حل نشده.

\subsection{یک‌شنبه ۳۰ دی}
سر یک \cl
\LTRfootnote{\url{https://go-review.googlesource.com/c/build/+/643415}}
 یکم اذیت کرد دوست عزیزمون \lr{Dmitri Shuralyov}. قبلا هم راجع بهش صحبت کرده بودیم. آدم خیلی خوبی هست با اینکه گیر زیاد میده
 ولی خیلی میشه ازش نکات زیبایی یاد گرفت. داخل کامنت ها برید ببینید دوتا مفهوم قشنگ رو اشاره کرده. یکیش \lr{information density} هست.
 و اینکه اصلا چی رو فدای چی باید کرد و چی چقدر مهمه. آخرین تغییراتم منجر به مشکل تو بیلد شده و روی یک پلتفرم خاص ارور داده. البته با این وضع
 هم این آقا امتیاز لازم برای مرج شدن رو دادن. خب خود این داداشمون متوجه شده که خطا در بیلد تو اون پلتفرم ربطی به من نداره و تائید کرده.
 این هم چند روز دیگه تائید میشه تهش. باید برم سراغ فیکس های بهتر دیگه. از فردا عصر شروع میکنم.

\subsection{دوشنبه ۱ بهمن}
خب امروز زمان کمی برای پیگیری مشارکت متن باز دارم. خطا بیلد رو \lr{Dmitri} امروز اوکی کرد و کامنت گذاشت ربطی به \cl من نداره.
البته مرج شدن این یکی خیلی طول کشیده واقعا. حالا چیز مهمی هم نیست. برای تغییرات دیگه داخل خود سورس کد گو تست ها مدنظرم هستن.
میتونیم با دستور \lr{go test -cover} کاوریج تست بگیریم و متوجه بشیم چه بخش هایی از یه کد تست نشده. شاید بشه اینطوری تو تست ها مشارکت کرد.
به‌علاوه که الان بر اساس رلیز سایکل داخل تایم فریز هستیم و فیچر نمیشه اضافه کرد به مخازن اصلی و حساس تر. بهترین کار الان تست نوشتن
و کار کردن روی داکیومنت ها هست. من با تست شروع می‌کنم و سعی می‌کنم \lr{coverage test} رو بیشتر کنم.

میدونیم که برای اجرا کردن تست ها یا هرکار دیگه ای روی کدبیس مخزن \lr{go} نیاز به همون ورژن کامپایلر داریم.
خب کافیه مخزن رو کلون کنید، کامپایل کنید و از باینری که دارید استفاده کنید. من مثلا خودم رو میزنم اینجا:
\begin{latin}
\begin{minted}{bash}
git clone https://go.googlesource.com/go
cd go/src
./make.bash
GO=$(pwd)/../bin/go
$GO test -cover ./... -coverprofile=/tmp/out && $GO tool cover -html=/tmp/out
$GO vet ./...
\end{minted}
\end{latin}

داستان کلا از این قرار هست. پس رسما با کامپایل کردن سورس کد گو البته با ورژن کامپایلر گو ای که درحال حاضر دارید، میتونید تست ها رو ران کنید
و از کدبیس استفاده کنید کلا.

\subsection{سه‌شنبه ۲ بهمن}
خب امروز تایم زیادی رو با درد سپری کردم و دیشب رسما از گلودرد نخوابیدم. بالاخره ظهر رفتم دکتر و تا ۲ اونجا بودم. بعدش هم
دوتا آنپول و قرص دوا حالم رو بهتر کرد.

همون \cl دیروز رو درگیرش بودم. یک نکته خیلی جالب که وقتی داشتم کتاب \lr{Go Design Pattern} رو مطالعه می‌کردم اسم
دونفر به‌عنوان کسایی که کامپایلر گو رو توسعه دادن و هد تیم بودن قید شد. یکی \lr{Ian Lance Taylor} بود و دیگری \lr{Russ Cox}.
جالبه که این دو نفر بعد از سالها هنوز خیلی فعال هستن و هد توسعه گولنگ هستن. جالبه که تا الان همه \cl های من رو همین \lr{Ian} مرج کرده.
و آخرین \cl من رو هم هردوشون تو \lr{reviewer} ها هستن و واقعا شگفت انگیزه این موضوع. همین \lr{Ian} امروز یک چیز جالب بهم یاد داد.
خب من تست رو نوشتم رو فرستادم. اگر برید داخل \cl مشخصه کامنت ها و جواب ها ولی اینجا یک مقدار توضیح میدم.

یک تابع ساده بود به اسم \lr{unhex} که داخل یک تابع دیگه به اسم \lr{unescape} بود استفاده شده بود و این یکی داخل
چند تابع دیگه مثل \lr{PathUnescape} استفاده شده بود. وقتی تست رو فرستادم، یان گفت که هدف ما از تست نوشتن رسیدن به پوشش ۱۰۰٪ نیست.
همین یک نکته. و بعد هم گفت بهتره تست توابعی که داخلشون استفاده شده رو اصلاح کنی. اما یک نکته ای که بود این بود که اون کیس خطا
هیچ وقت داخل تابع \lr{unescape} اتفاق نمیوفتاد. چرا که قبل از رسیدن به این تابع \lr{validation} داشتیم.
در این مورد یان گفت اگر همچین وضعیتی داریم، میتونی کیس خطا رو با \lr{panic} عوض کنی. دلیلش هم این هست که هیچ وقت رخ نمیده و
در صورتی که اتفاق بیوفته هم عملکرد پیشبینی نشده هست.

خلاصه امشب یان تائید رو داد. حالا باید صبر کنم یکی دیگه از کارمندان گوگل کارم رو تائید کنه. امروز روز کمکاری بود اما بیکار نبودم.
کار اوپن سورس رو فعلا واسه امروز در همین حد کافی میدونم و میرم کتابم رو به اندازه کافی بخونم. حداقل ۴۰ صفحه دیگه باید تموم بشه.
فردا صبح زود هم انتخاب واحد دارم.


